\documentclass[11pt,a4paper]{article}
\usepackage[margin=1in]{geometry}
\usepackage{hyperref}
\usepackage{enumitem}

\title{Strategic Management \\ Week 10 \\ Slide-by-Slide Notes (ELI5) \\ Unit 2}
\author{(Generated from \texttt{sm-week-10-slides-3.pdf})}
\date{}

\begin{document}
\maketitle

\section*{How to use this (ELI5)}
\begin{itemize}
  \item For each slide, write the ``Exam write'' line in your notebook.
  \item Then copy the ELI5 explanation in simple English.
  \item If you get stuck, re-read only that slide section.
  \item Aim for: label + reason + exam sentence.
\end{itemize}

% ---------------------------
\section*{Slide 1 of 14 (Contents: Unit 2)}
\textbf{Key idea (from slide):} Unit 2 is Objectives and Values.\\
\textbf{ELI5:} This unit is about why companies have goals, and how values affect decisions.\\
\textbf{Exam write:} ``Unit 2 covers mission/vision/objectives, value creation with stakeholders, governance, and corporate values.''

% ---------------------------
\section*{Slide 2 of 14 (Learning Objectives)}
\textbf{Key idea (from slide):}
\begin{itemize}
  \item economic and social perspectives (shareholders vs stakeholders)
  \item how to calculate economic profit
  \item mechanisms of corporate governance
  \item integrate sustainability and CSR into mission/vision
\end{itemize}
\textbf{ELI5:} You learn what companies try to optimize, and how they control decisions.\\
\textbf{Exam write:} ``I will connect objectives and values to governance and to stakeholder vs shareholder perspectives.''

% ---------------------------
\section*{Slide 3 of 14 (2.1 Mission, Vision and Strategic Objectives)}
\textbf{Key idea (from slide):} Corporate mission variables: scope, core capabilities, values/beliefs/behavior; mission is stable over time.\\
\textbf{ELI5:} Mission is ``who we are'' and ``what we do'' right now.\\
\textbf{Exam write:} ``Mission = present identity and business essence, kept stable over time.''

% ---------------------------
\section*{Slide 4 of 14 (Corporate vision)}
\textbf{Key idea (from slide):} Corporate vision = what the firm will be in 10 years+; CEO main task is visionaire; must be clear and inspiring and inclusive.\\
\textbf{ELI5:} Vision is the future dream, and it guides teamwork.\\
\textbf{Exam write:} ``Vision is the long-run destination that must be clear and motivating.''

% ---------------------------
\section*{Slide 5 of 14 (Exhibit 2.1 mission examples)}
\textbf{Key idea (from slide):} Exhibit showing examples of corporate mission statements.\\
\textbf{ELI5:} Examples help you see what ``good mission'' looks like.\\
\textbf{Exam write:} ``Good missions are stable and explain the essence of the business.''

% ---------------------------
\section*{Slide 6 of 14 (Strategic objectives)}
\textbf{Key idea (from slide):} Objectives explain how we close the gap from present to future; they have: measurable attribute, yardstick/metric, target value, timeframe.\\
\textbf{ELI5:} Objectives are the small steps with numbers and deadlines.\\
\textbf{Exam write:} ``Objectives must be measurable, have a metric, a target level, and a time frame.''

% ---------------------------
\section*{Slide 7 of 14 (Value creation and stakeholders: Firm performance importance + issues)}
\textbf{Key idea (from slide):}
\begin{itemize}
  \item performance measurement matters to judge success/failure and CEO/TMT performance
  \item difficult to measure and unclear who defines performance
  \item success depends on how performance is defined
\end{itemize}
\textbf{ELI5:} The hardest part is choosing ``success'' and measuring it.\\
\textbf{Exam write:} ``In answers, always say what performance definition you use (and why it matters).''

% ---------------------------
\section*{Slide 8 of 14 (Performance measurement: accounting, market, social/environment)}
\textbf{Key idea (from slide):}
\begin{itemize}
  \item accounting profitability: gross operating margin, EBIT, net income, ROA, ROE
  \item firm value: market capitalization and shareholder value creation (share price)
  \item also social/environment: positive impact, shared value, minimize pollution emissions
\end{itemize}
\textbf{ELI5:} Success can be money numbers OR social/environment numbers, depending on the viewpoint.\\
\textbf{Exam write:} ``Performance can be measured as accounting profitability, market value, and social/environment impact.''

% ---------------------------
\section*{Slide 9 of 14 (Stakeholders + shareholder vs stakeholder theories)}
\textbf{Key idea (from slide):}
\begin{itemize}
  \item stakeholders are people/groups affecting or affected by the firm
  \item Chicago School (Milton Friedman): managers maximize shareholder wealth
  \item Stakeholders’ Theory (Freeman): triple bottom line (economic, social, environmental), stakeholder satisfaction, long-term orientation
\end{itemize}
\textbf{ELI5:} Shareholder theory says ``maximize shareholder wealth''.\\
\textbf{ELI5:} Stakeholder theory says ``serve everyone who matters'': workers, customers, society.\\
\textbf{Exam write:} ``I will contrast shareholder primacy vs stakeholder theory and explain the objective implication.''

% ---------------------------
\section*{Slide 10 of 14 (Stakeholders: Illustration / critical thinking role-play)}
\textbf{Key idea (from slide):} Role play: make positive and negative arguments about high corporate social responsibility.\\
\textbf{ELI5:} CSR is not only good or only bad; you should argue both sides.\\
\textbf{Exam write:} ``When asked about CSR, I will discuss tradeoffs: purpose vs profit, and stakeholder impacts.''

% ---------------------------
\section*{Slide 11 of 14 (2.3 Corporate Governance: ownership separation and mechanisms)}
\textbf{Key idea (from slide):}
\begin{itemize}
  \item ownership -> separation -> management
  \item governance mechanisms: incentives (bonuses, stock options), board of directors supervision, gate between owners and managers
  \item conflict: who sets the objectives?
  \item shareholder value creation linked to remuneration vs power/status
\end{itemize}
\textbf{ELI5:} Governance is the ``control system'' that makes managers think like owners.\\
\textbf{Exam write:} ``Corporate governance reduces the owner-manager incentive conflict through incentives and board supervision.''

% ---------------------------
\section*{Slide 12 of 14 (2.4 Corporate Values: what they do)}
\textbf{Key idea (from slide):} Corporate values guide business activities, ethical behavior, stakeholder engagement, and organizational culture; ``the end does not justify the means''.\\
\textbf{ELI5:} Values are the rules for how the company acts, not only what it wants.\\
\textbf{Exam write:} ``Corporate values guide ethical behavior and stakeholder engagement and shape culture.''

% ---------------------------
\section*{Slide 13 of 14 (Business ethics + CSR + move from profit maximization to purpose)}
\textbf{Key idea (from slide):}
\begin{itemize}
  \item voluntary activities to meet social/environment demands
  \item move from just profit maximization -> purpose
  \item move from shareholder primacy -> stakeholder engagement (CSR)
  \item business ethics: what is right/wrong in strategic decisions and transactions
\end{itemize}
\textbf{ELI5:} CSR and ethics tell the company how to behave, not only how to make profit.\\
\textbf{Exam write:} ``CSR integrates purpose and stakeholder engagement, and business ethics defines right vs wrong behavior in decisions.''

% ---------------------------
\section*{Slide 14 of 14 (Business case: Danone + questions)}
\textbf{Key idea (from slide):} Business case Danone: conflict between profit and purpose; board defense of strategy pursuing purpose.\\
\textbf{ELI5:} You must argue: how can a board justify purpose-oriented strategy?\\
\textbf{Exam write:} ``In the Danone-style question: explain profit vs purpose conflict, then propose governance logic to defend purpose strategy.''

\end{document}

